\documentclass[a4paper,10pt]{article}
\usepackage[utf8]{inputenc}
\usepackage{graphicx}
\usepackage{subfigure}
\usepackage{fancyhdr}
\usepackage{listings}
\usepackage{ gensymb }


%opening

\pagestyle{fancy} %eigener Seitenstil
\fancyhf{} %alle Kopf- und Fußzeilenfelder bereinigen
\fancyhead[L]{SEP Abgabegespräch SS2015/16} %Kopfzeile links
\fancyhead[C]{} %zentrierte Kopfzeile
\fancyhead[R]{\includegraphics[width=1cm]{../../template/tugraz_bw}} %Kopfzeile rechts
\renewcommand{\headrulewidth}{0.4pt} %obere Trennlinie
\fancyfoot[C]{} %Seitennummer
\renewcommand{\footrulewidth}{0.4pt} %untere Trennlinie


\lstset{language=C++, basicstyle=\ttfamily\tiny}  

\begin{document}

\section*{CSV Parser}

Entwicklen Sie eine Software die eine csv Datei einliest. 
CSV Dateien sind dadurch gekennzeichnet, dass die einzelnen Werte (Spalten) in einer Zeile
durch ein Semicolon voneinander getrennt sind. Konkreter, nach jedem Spaltenwert (kann auch leer sein) steht ein Semicolon.
Also auch am Ende der Zeile. Ob alle Zeilen die gleiche Spaltenanzahl haben wird nicht getestet. Sie können davon ausgehen.

Sollte es zu einer Fehlermeldung bei der Abarbeitung eines Befehles kommen
wird dieser "abgebrochen" und es wird wieder auf eine Eingabe gewartet.
Sollte ein Befehl nicht dem vorgegebenen Format entsprechen wird er einfach ignoriert. Case Sensitive Unterstützung bei den Commandos (nicht den Argumenten) ist optional.
Nachfolgend die zu implementierenden Befehle:

\begin{description}
 \item[load $<$csvfile$>$] Lädt das angegebene CSV File. Sollte dieses File nicht existieren wird "File does not exist!" ausgegeben.
 \item[show $<$linenumber$>$ $[$ $<$columnnumber$>$]$ $] Zeigt entweder die ganze Zeile an oder wenn auch die Spaltennummer angegeben wurde nur den jeweiligen Spalteninhalt.
 Sollte eine der Eingaben nicht gültig sein (weil sie nicht existieren) wird "Numbers not valid!" ausgegeben. Spalten und Zeilenangaben beginnen bei 1.
\end{description}

Orientieren Sie sich zusätzlich an dem unten vorgegebenen Beispiel. Bei der Eingabe von \texttt{quit} oder EOF soll das Programm beendet werden.

\subsection*{valid.csv:}
\begin{verbatim}
SEP;AGS;SIND;SUPER;
CSV;PARSER;SIND;WICHTIG;
\end{verbatim}

\subsection*{Beispiel}
\begin{verbatim}
> load file_doest_not_exist.csv
File does not exist!
> load valid.csv
> show 2
CSV;PARSER;SIND;WICHTIG;
> show 2 2
PARSER
\end{verbatim}

\subsection*{Aufruf}
\texttt{./csvParser}

\subsection*{Bewertung}
\begin{itemize}
 \item Dokumentation und Programmierstil: 1 Punkt
 \item Einhalten der OOP Konzepte: 1 Punkt
 \item Struktur \& Korrektheit: 1 Punkt
\end{itemize}

\newpage


\end{document}
